% =====================================================================
%  Solution template -- Nestruev, "Smooth Manifolds and Observables"
%  Differential Geometry Reading Group
%
%  1. Save as solutions/chNN/<exercise>_<your-name>.tex
%     e.g. solutions/ch03/3-28_aleksandras.tex
%  2. Fill in the part marked YOUR SOLUTION at the bottom.
%  The preamble is shared by everyone: please don't change it here.
%  If you need new notation, add it under "Your own macros".
% =====================================================================
\documentclass[11pt,a4paper]{article}

% ------------------------------------------------ shared preamble (don't edit)
\usepackage[margin=2.5cm]{geometry}
\usepackage{amsmath,amssymb,amsthm,mathtools}
\usepackage{enumitem}
\usepackage{tikz-cd}
\usepackage{xcolor}
\usepackage[colorlinks=true,linkcolor=blue!50!black,urlcolor=blue!50!black]{hyperref}
\usepackage[capitalise]{cleveref}

% ---------------------------------------------------------------- metadata
\makeatletter
\newcommand{\dg@chapnum}{}
\newcommand{\dg@chaptitle}{}
\newcommand{\chapter}[2]{\renewcommand{\dg@chapnum}{#1}\renewcommand{\dg@chaptitle}{#2}}
\renewcommand{\maketitle}{%
  \begin{center}
    {\small Nestruev, \emph{Smooth Manifolds and Observables} --- reading group}\par
    \vspace{0.6em}
    {\Large\bfseries Chapter \dg@chapnum: \dg@chaptitle}\par
    \vspace{0.4em}
    {\normalsize \@author}\par
    {\small \@date}\par
  \end{center}
  \vspace{1em}}
\makeatother

% ---------------------------------------------------------------- environments
\theoremstyle{plain}
\newtheorem{theorem}{Theorem}
\newtheorem{lemma}[theorem]{Lemma}
\newtheorem{proposition}[theorem]{Proposition}
\newtheorem{corollary}[theorem]{Corollary}
\newtheorem{claim}[theorem]{Claim}
\theoremstyle{definition}
\newtheorem{definition}[theorem]{Definition}
\newtheorem{example}[theorem]{Example}
\theoremstyle{remark}
\newtheorem*{remark}{Remark}

% \begin{exercise}{3.28} ... \end{exercise}   -- restate the exercise from the book
\newenvironment{exercise}[1]
  {\par\medskip\noindent\textbf{Exercise #1.}\ \itshape}
  {\par\medskip}
% \begin{solution} ... \end{solution}
\newenvironment{solution}
  {\begin{proof}[Solution]}
  {\end{proof}}

% ---------------------------------------------------------------- notation
% Notation roughly follows Nestruev. Add to this list rather than redefining
% things in individual solution files.
\newcommand{\R}{\mathbb{R}}
\newcommand{\C}{\mathbb{C}}
\newcommand{\Z}{\mathbb{Z}}
\newcommand{\N}{\mathbb{N}}
\newcommand{\Q}{\mathbb{Q}}
\newcommand{\RP}{\mathbb{R}P}
\newcommand{\F}{\mathcal{F}}            % algebra of observables
\newcommand{\Cinf}{C^{\infty}}
\newcommand{\abs}[1]{\left\lvert #1 \right\rvert}
\newcommand{\norm}[1]{\left\lVert #1 \right\rVert}
\newcommand{\Spec}{\operatorname{Spec}}
\newcommand{\Specm}{\operatorname{Specm}}
\newcommand{\Hom}{\operatorname{Hom}}
\newcommand{\End}{\operatorname{End}}
\newcommand{\Der}{\operatorname{D}}     % derivations D(P)
\newcommand{\Diff}{\operatorname{Diff}} % differential operators Diff_k(P,Q)
\newcommand{\smbl}{\operatorname{smbl}}
\newcommand{\supp}{\operatorname{supp}}
\newcommand{\id}{\mathrm{id}}
\newcommand{\im}{\operatorname{im}}
\newcommand{\coker}{\operatorname{coker}}
\newcommand{\Ann}{\operatorname{Ann}}
\newcommand{\dd}{\,\mathrm{d}}
\newcommand{\pd}[2]{\frac{\partial #1}{\partial #2}}
% |F| : the R-spectrum (dual space) of an R-algebra F
\newcommand{\dual}[1]{\lvert #1 \rvert}

% ------------------------------------------------ your own macros (optional)


% =====================================================================
%  YOUR SOLUTION
% =====================================================================
\chapter{1}{Introduction}
\author{Aleksandras}
\date{\today}

\begin{document}
\maketitle

\begin{exercise}{1.2}
In nine-dimensional Euclidean space $\mathbb{R}^9$ consider the set of points
satisfying the following system of six algebraic equations:
\[
\begin{cases}
x_1^2 + x_2^2 + x_3^2 = 1; & x_1x_4 + x_2x_5 + x_3x_6 = 0;\\
x_4^2 + x_5^2 + x_6^2 = 1; & x_1x_7 + x_2x_8 + x_3x_9 = 0;\\
x_7^2 + x_8^2 + x_9^2 = 1; & x_4x_7 + x_5x_8 + x_6x_9 = 0.
\end{cases}
\]
 describe a bijective map of a neighborhood of any
point (say $(1, 0, 0, 0, 1, 0, 0, 0, 1)$) of the surface onto a neighborhood of the
origin of Euclidean 3-space. But this map cannot be extended to cover the
entire surface, which is compact (why?)
\end{exercise}

\begin{solution}
  Key idea is that we first interpret $(x_1, x_2, x_3)$ as a vector in $\R^3$, i.e. we set $A=(x_A, y_A, z_A):=(x_1, x_2, x_3)$. The same goes for the other variables:
  \begin{itemize}
    \item $B=(x_B, y_B, z_B):=(x_4, x_5, x_6)$
    \item  $C=(x_C, y_C, z_C):=(x_7, x_8, x_9)$
  \end{itemize}

  Then the equations tell us the all three vectors $A,B$ and $C$ lie on the unit sphere around the origin and that they are pairwise orthogonal. This essentialy brings us to the example 1.1 of the rotating body so we really only need two parameters to say how the first vector is oriented and one more parameter to say how the second vector is oriented around the axis of the first one. While fixing $A$ and $B$ leaves 2 choices for $C$ the excercise asks us to make it a mapping from the neighbourhood of the 9 dimensional point  $p=(1, 0, 0, 0, 1, 0, 0, 0, 1)$ so only one of those choices for $C$ will be close to $p$.

  It should be noted that under the mapping in example 1.1. the 9-dimesional vector gets sent to $(x_1, x_2, \phi)$ which are close to $(0,0,0)$ for a small deviation from $p$, thus this is indeed a map to the neighbourhood of $(0,0,0)\in \R$

  The mapping fails to extend to the entire space exactly for the same reasons that it does for example 1.1.

  \end{solution}
  \end{document}
